Táto príloha uvádza prehľad štruktúry elektronického nosiča odovzdávaného
riešenia. Zdrojové kódy systému FinanceGPT sú organizované ako projekt webovej aplikácie
postavenej na frameworku Ruby on Rails a rozšírenej o komponenty pre spracovanie dotazov
v prirodzenom jazyku, chatové rozhranie a analytické nástroje nad finančnými
dátami.

\subsection*{A.1 Hlavná štruktúra projektu}

\begin{verbatim}
/app
  /controllers      - kontroléry aplikácie, spracovanie HTTP požiadaviek
  /models           - dátové modely aplikácie
  /services         - servisné triedy pre AI a analytické nástroje
  /views            - šablóny používateľského rozhrania

/app/views/chat
  index.html.erb    - hlavné konverzačné rozhranie systému FinanceGPT

/javascript
  /controllers      - Stimulus kontroléry pre frontendové správanie

/config
  routes.rb         - definícia aplikačných trás
  initializers/     - inicializačné súbory konfigurácie

/db
  /migrate          - databázové migrácie
  seeds.rb          - inicializačný skript s 
                      testovacími dátami (100 faktúr, 12 klientov)

/lib
  /tasks            - pomocné Rake úlohy

/test
  - jednotkové a integračné testy

README.md           - základný popis projektu a pokyny na spustenie
Gemfile             - definícia Ruby závislostí projektu
\end{verbatim}

\subsection*{A.2 Hlavné aplikačné komponenty}

V adresári \texttt{app/controllers} sa nachádzajú kontroléry zabezpečujúce tok
požiadaviek medzi používateľským rozhraním, aplikačnou logikou a databázovou vrstvou.
Ide najmä o kontroléry pre domovskú stránku, používateľov, faktúry a chatové
rozhranie.

Adresár \texttt{app/models} obsahuje modely reprezentujúce základné doménové objekty
systému. Patria sem modely používateľa, faktúry, bankových údajov, správ v konverzácii,
samostatných chatov a záznamov o volaniach nástrojov.

Adresár \texttt{app/services} obsahuje servisné triedy implementujúce logiku spracovania
používateľských dotazov. Nachádzajú sa tu nástroje zamerané na generovanie SQL
dopytov, sémantické vyhľadávanie, rozklad príjmov a generovanie grafických
výstupov cashflow.

V adresári \texttt{app/views} sa nachádzajú šablóny používateľského rozhrania.
Osobitne dôležitou časťou je priečinok \texttt{app/views/chat}, ktorý obsahuje
šablónu hlavného chatového rozhrania inteligentného asistenta.

Frontendová interaktivita je sústredená v adresári \texttt{javascript/controllers},
kde sa nachádzajú Stimulus kontroléry. Táto vrstva zabezpečuje dynamické správanie
chatového rozhrania a nadväzuje na architektúru Hotwire/Turbo použitú v aplikácii.

Databázová vrstva je reprezentovaná adresárom \texttt{db/migrate}, ktorý obsahuje
migračné súbory definujúce vývoj databázovej schémy. Súbor \texttt{seeds.rb} slúži
na inicializáciu vybraných dát potrebných na testovanie alebo demonštráciu systému.

Adresár \texttt{lib/tasks} obsahuje pomocné úlohy spúšťané prostredníctvom
nástroja Rake. V rámci projektu sem patria aj úlohy súvisiace so spracovaním embeddingov
a prípravou dát pre sémantické vyhľadávanie.

Adresár \texttt{test} obsahuje testovacie scenáre pre vybrané časti aplikácie. Testovanie
zahŕňa jednotkové aj integračné testy zamerané na korektnosť spracovania dát a funkčnosť
jednotlivých komponentov systému.

\subsection*{A.3 Obsah odovzdávaného riešenia}

Elektronický nosič obsahuje úplné zdrojové kódy aplikácie, konfiguračné súbory,
databázové migrácie, pomocné skripty a testy. Súčasťou odovzdaného riešenia je
aj súbor \texttt{README.md}, ktorý slúži ako stručný technický sprievodca pre
lokálne spustenie projektu a orientáciu v jeho štruktúre.